\section*{Hoofdstuk \ref{ch:b3:complex-spectra-magnetism} --- Elektronische spectra en magnetisme van complexen}

\begin{solution}{exo:b3:complex-spectra-magnetism:1}
D: E$_g$ + T$_{2g}$ ($2 + 3 = 5$); F: A$_{2g}$ + T$_{1g}$ + T$_{2g}$ ($1 + 3 + 3 = 7$); G: A$_{1g}$ + E$_g$ + T$_{1g}$ + T$_{2g}$
($1 + 2 + 3 + 3 = 9$).
\end{solution}

\begin{solution}{exo:b3:complex-spectra-magnetism:2}
$d^1$ $^2$D $\to$ $^2$T$_{2g}$; $d^2$ $^3$F $\to$ $^3$T$_{1g}$; $d^3$ $^4$F $\to$ $^4$A$_{2g}$; $d^4$ $^5$D $\to$ $^5$E$_g$; $d^5$ $^6$S $\to$ $^6$A$_{1g}$; $d^6$ $^5$D
$\to$ $^5$T$_{2g}$; $d^7$ $^4$F $\to$ $^4$T$_{1g}$; $d^8$ $^3$F $\to$ $^3$A$_{2g}$; $d^9$ $^2$D $\to$ $^2$E$_g$.
\end{solution}

\begin{solution}{exo:b3:complex-spectra-magnetism:3}
Hoogspin-\ce{Mn^{2+}} ($n = 5$) $5.92\,\mu_B$, \ce{Fe^{2+}} (4) $4.90\,\mu_B$, \ce{Co^{2+}} (3) $3.87\,\mu_B$; laagspin-\ce{Fe^{2+}} en
-\ce{Co^{3+}} ($t_{2g}^6$): 0, diamagnetisch.
\end{solution}

\begin{solution}{exo:b3:complex-spectra-magnetism:4}
$\ce{[Mn(H2O)6]^{2+}}$ (spinverboden en verboden door Laporte, bijna kleurloos) $<$ $\ce{[Co(H2O)6]^{2+}}$
(verboden door Laporte, vibronisch) $<$ $\ce{[CoCl4]^{2-}}$ (geen symmetriecentrum) $<$ \ce{MnO4-} (toegestane
ladingsoverdracht).
\end{solution}

\begin{solution}{exo:b3:complex-spectra-magnetism:5}
8500: $^3$A$_{2g}\to{}^3$T$_{2g}$, dus $\Delta_{\mathrm o} = \qty{8500}{cm^{-1}}$; 14100: $^3$T$_{1g}$(F); 24900: $^3$T$_{1g}$(P).
$15B = 39\,000 - 3 \times 8500$, $B = \qty{900}{cm^{-1}}$, $\beta = 900/1056 = 0.85$.
\end{solution}

\begin{solution}{exo:b3:complex-spectra-magnetism:6}
$\Delta_{\mathrm o} = \qty{17400}{cm^{-1}}$; oplossen van $7.5B + 1.5\Delta_{\mathrm o} - \frac12\sqrt{225B^2 - 18B\Delta_{\mathrm o} + \Delta_{\mathrm o}^2} = 24\,600$ geeft $B = \qty{729}{cm^{-1}}$,
$\beta = 0.79$. $\nu_3 = \qty{38500}{cm^{-1}}$ (\qty{260}{nm}), in het ultraviolet.
\end{solution}

\begin{solution}{exo:b3:complex-spectra-magnetism:7}
De $d$-elektronen zijn sterker gedelokaliseerd naar de oxide-ionen van de
edelstenen dan naar watermoleculen: de bindingen Cr--O zijn er meer
covalent.
\end{solution}

\begin{solution}{exo:b3:complex-spectra-magnetism:8}
Sterk ($e_g$ ongelijk gevuld): $d^9$, hoogspin-$d^4$, laagspin-$d^7$. Geen: $d^3$ ($t_{2g}^3$, term A) en
laagspin-$d^6$ ($t_{2g}^6$). Zwak: hoogspin-$d^6$ ($t_{2g}^4e_g^2$, term T).
\end{solution}

\begin{solution}{exo:b3:complex-spectra-magnetism:9}
Octaëdrisch hoogspin-$d^7$: grondterm $^4$T$_{1g}$, een term T met een \omterm{def:b3:complex-spectra-magnetism:moment}{orbitale bijdrage},
een moment ruim boven de spin-only-waarde $3.87\,\mu_B$. Tetraëdrisch $d^7$ ($\equiv$ octaëdrisch $d^3$):
grondterm $^4$A$_2$, orbitaal moment uitgedoofd, dicht bij de spin-only-waarde
(iets erboven, door menging met aangeslagen termen T).
\end{solution}

\begin{solution}{exo:b3:complex-spectra-magnetism:10}
$\mu_{\mathrm{eff}} = 2.828\sqrt{1.87} = 3.87\,\mu_B$: drie ongepaarde elektronen ($S = \frac32$).
\end{solution}

\begin{solution}{exo:b3:complex-spectra-magnetism:11}
$\Delta\chi_v = 3 \times 25.0/\num{400e6} = \num{1.88e-7}$; $\chi_{\mathrm m} = \num{1.88e-7}/\qty{10.0}{mol.m^{-3}} = \qty{1.88e-8}{m^3.mol^{-1}}$, d.w.z.\
\qty{1.49e-3}{cm^3.mol^{-1}}; $\chi_{\mathrm m}T = 0.445$, $\mu_{\mathrm{eff}} = 1.89\,\mu_B$: één ongepaard elektron.
\end{solution}

\begin{solution}{exo:b3:complex-spectra-magnetism:12}
$T_{1/2} = 15\,000/75 = \qty{200}{K}$. Bij \qty{250}{K}: $\Delta G = 15\,000 - 250 \times 75 = \qty{-3750}{J/mol}$, $K = \eu^{1.80} = 6.07$,
$x_{\mathrm{HS}} = 0.86$; $\chi_{\mathrm m}T = 0.86 \times 3.00 = \qty{2.6}{cm^3.K.mol^{-1}}$.
\end{solution}

\begin{solution}{pb:b3:complex-spectra-magnetism:1}
\textbf{1.} [Ar]$3d^3$.
\textbf{2.} $^4$F.
\textbf{3.} $^4$A$_{2g}$ + $^4$T$_{2g}$ + $^4$T$_{1g}$.
\textbf{4.} $^4$A$_{2g}$, de term van $t_{2g}^3$, met alle drie de elektronen in de lagere orbitalen.
\textbf{5.} $^4$T$_{1g}$(P).
\textbf{6.} $^4$A$_{2g}\to{}^4$T$_{2g}$, $\to{}^4$T$_{1g}$(F), $\to{}^4$T$_{1g}$(P).
\textbf{7.} 561 en \qty{414}{nm}.
\textbf{8.} 597 en \qty{434}{nm}.
\textbf{9.} Robijn absorbeert geelgroen en violet en laat rood door (met wat
blauw); de banden van smaragd, naar het rood verschoven, absorberen ook
oranjerood, en groen wordt doorgelaten.
\textbf{10.} Robijn \qty{17820}{cm^{-1}}, smaragd \qty{16740}{cm^{-1}}.
\textbf{11.} $^4$T$_{2g}$ ($t_{2g}^2e_g$) en $^4$A$_{2g}$ ($t_{2g}^3$) hebben dezelfde elektronenafstotingsenergie; hun
verschil is de promotie-energie van één elektron, $\Delta_{\mathrm o}$.
\textbf{12.} $B = (14\,305 - 547)/15 = \qty{917}{cm^{-1}}$.
\textbf{13.} \qty{614}{cm^{-1}}.
\textbf{14.} \qty{618}{cm^{-1}}.
\textbf{15.} 0.67 voor beide.
\textbf{16.} 29.0 en 27.1.
\textbf{17.} $\nu_3 = \qty{38500}{cm^{-1}}$ (\qty{260}{nm}) voor robijn, in het ultraviolet, waar de
ladingsoverdrachtsabsorptie hem verbergt.
\textbf{18.} De octaëdrische plaats van beryl is groter: langere bindingen
Cr--O geven een kleinere splitsing, die steil varieert met de afstand.
\textbf{19.} $10^7/695 = \qty{14390}{cm^{-1}}$, $23.4B$.
\textbf{20.} $^2$E$_g$ hangt vooral af van $C$: het hoger gemeten niveau betekent dat $C/B$ groter is,
ongeveer 5.3 volgens dezelfde diagonalisatie.
\textbf{21.} $^2$E$_g$ en $^4$A$_{2g}$ behoren allebei tot $t_{2g}^3$ en verschillen alleen door het omklappen van een spin: de binding,
en dus de geometrie, is in beide dezelfde, en de overgang heeft geen
vibrationele progressie (een 0--0-lijn).
\textbf{22.} Ze is spinverboden.
\textbf{23.} Haar energie hangt af van $B$ en $C$, bijna niet van $\Delta_{\mathrm o}$ (een vlakke lijn in het diagram),
en $B$ is in de twee edelstenen bijna gelijk.
\textbf{24.} Ionen die in de brede banden gepompt worden, relaxeren en hopen
zich op in het langlevende
niveau $^2$E$_g$, zodat er meer ionen in kunnen zitten dan in het grondniveau: een
populatie-inversie.
\textbf{25.} \textbf{$\Delta_{\mathrm o}(\text{robijn}) - \Delta_{\mathrm o}(\text{smaragd}) = \qty{1080}{cm^{-1}}$}: het hele verschil tussen rood en
groen.
\end{solution}
